\subsection*{一步一步建立直观}

\paragraph{1. 先选一个可重复测量的面积单位}

观察单位方形的四个顶点在同一矩阵下的去向。输出是由两列张成的平行四边形，它与原方形的面积之比不依赖方形放在哪里。行列式正是把这种统一的局部面积倍率压缩成一个数。



\paragraph{2. 剪切说明形状改变未必改变面积}

剪切把顶部推向一侧，却保留底边和垂直高度，因此面积仍为一。只看两根边的长度会误判，因为夹角同时改变。坐标公式 $ad-bc$ 把两块有向矩形相减，恰好消除了这种倾斜造成的重复计算。



\paragraph{3. 用反射区分大小与定向}

反射后方形面积仍为一，但两根边的顺时针或逆时针次序交换，因此行列式是负一。普通面积必须取绝对值。若继续把两列调成共线，面积为零；这又把几何退化与不能唯一恢复输入联系起来。
